\subsection{RISC-V 等价实现（成员A负责）}
\subsubsection{SysY 到 RISC-V 的实现对应}
RISC-V 版本保持与 \texttt{main.sy} 相同的函数划分和控制逻辑。按照 RV64 调用约定，整数和地址参数通过 \texttt{a0}、\texttt{a1} 等参数寄存器传递，整数返回值放在 \texttt{a0} 中；\texttt{main} 中需要跨函数调用保存的 $n$、循环变量和数组基址分别放入 \texttt{s0}、\texttt{s1}、\texttt{s2}，并在栈上保存 \texttt{ra} 和这些被调用者保存寄存器。栈帧大小取 80 B，为 16 的整数倍。

\begin{table}[H]
\centering
\caption{SysY 结构与 RISC-V 实现的对应关系}
\begin{tabularx}{0.96\linewidth}{>{\ttfamily}p{3.1cm} >{\ttfamily}p{4.2cm} X}
\toprule
\normalfont SysY 结构 & \normalfont RISC-V 实现 & \normalfont 作用 \\
\midrule
函数参数 a、n & a0、a1 & 按照调用约定传递数组首地址和整数参数 \\
return & a0 & 函数返回值寄存器 \\
a[i] & slli + add + lw/sw & 将下标左移 2 位得到 $i\times4$，再与数组基地址相加 \\
while & 条件分支 + 标签 + 回跳 & 将循环转换为基本块和跳转 \\
a[i]\%2==0 & remw + bnez & 求余后判断是否为 0 \\
a[i]>0 & blez & 非正数时跳过计数 \\
getint/putint/\newline putch & call & 调用 SysY 运行时库函数 \\
main 栈帧 & addi sp,sp,-80 / +80 & 保存 ra 和 s0--s2，并保持 16 B 对齐 \\
\bottomrule
\end{tabularx}
\end{table}

汇编中使用了 \texttt{li}、\texttt{mv}、\texttt{j}、\texttt{call}、\texttt{ret}、\texttt{bnez}、\texttt{blez} 等常见伪指令，汇编器会将其展开为相应的基础指令组合。手写汇编的重点不是逐字翻译源代码，而是保持函数接口、控制流、数据访问和最终运行结果与 SysY 程序等价。

\subsubsection{RISC-V 汇编程序}
\begin{lstlisting}[language=RiscVAsm,caption={与 main.sy 等价的 RV64 汇编程序 main.s},label={lst:riscv-main}]
.option nopic
.attribute stack_align,16
.text

.globl sum_even
.type sum_even,@function
sum_even:
#初始化
    li t0,0
    li t1,0
.Lse_loop:
    bge t0,a1,.Lse_end
    slli t2,t0,2
    add t3,a0,t2
    lw t4,0(t3)
#判断偶数
    li t5,2
    remw t6,t4,t5
    bnez t6,.Lse_next
    addw t1,t1,t4
.Lse_next:
    addi t0,t0,1
    j .Lse_loop
.Lse_end:
    mv a0,t1
    ret
.size sum_even,.-sum_even

.globl count_pos
.type count_pos,@function
count_pos:
#初始化
    li t0,0
    li t1,0
.Lcp_loop:
    bge t0,a1,.Lcp_end
    slli t2,t0,2
    add t3,a0,t2
    lw t4,0(t3)
#统计正数
    blez t4,.Lcp_next
    addiw t1,t1,1
.Lcp_next:
    addi t0,t0,1
    j .Lcp_loop
.Lcp_end:
    mv a0,t1
    ret
.size count_pos,.-count_pos

.globl main
.type main,@function
main:
#开辟栈空间
    addi sp,sp,-80
    sd ra,72(sp)
    sd s0,64(sp)
    sd s1,56(sp)
    sd s2,48(sp)

    mv s2,sp

#输入n
    call getint
    mv s0,a0
    li s1,0

#输入数组
.Linput:
    bge s1,s0,.Linput_end
    call getint
    slli t0,s1,2
    add t0,s2,t0
    sw a0,0(t0)
    addi s1,s1,1
    j .Linput

.Linput_end:
#求偶数和
    mv a0,s2
    mv a1,s0
    call sum_even
    call putint
    li a0,10
    call putch

#统计正数
    mv a0,s2
    mv a1,s0
    call count_pos
    call putint
    li a0,10
    call putch

    li a0,0
    ld s2,48(sp)
    ld s1,56(sp)
    ld s0,64(sp)
    ld ra,72(sp)
    addi sp,sp,80
    ret
.size main,.-main
\end{lstlisting}

\subsubsection{汇编与目标文件生成}
首先使用 RISC-V 交叉编译器的汇编阶段把 \texttt{main.s} 转换为可重定位目标文件 \texttt{main.o}：
\begin{lstlisting}[language=bash]
riscv64-unknown-elf-gcc main.s -c -o main.o -w
file main.o
\end{lstlisting}
\texttt{file} 输出中包含 \texttt{ELF 64-bit}、\texttt{relocatable} 和 \texttt{UCB RISC-V}，说明文本汇编已被编码为 64 位 RISC-V 目标文件，但此时仍处于可重定位状态，尚不能独立运行。

\subsubsection{链接 SysY 运行时库}
程序中的 \texttt{getint}、\texttt{putint} 和 \texttt{putch} 不是在 \texttt{main.s} 中实现的，因此需要链接课程提供的 SysY RISC-V 运行时库 \texttt{libsysy\_riscv.a}。实验采用静态链接和 \texttt{medany} 内存模型，并按课程环境把代码基址设置为 \texttt{0x90000000}：
\begin{lstlisting}[language=bash]
riscv64-unknown-elf-gcc main.o -o main \
  -L./lib -lsysy_riscv \
  -static -mcmodel=medany \
  -Wl,--no-relax,-Ttext=0x90000000
\end{lstlisting}

首次尝试链接时，终端优先使用了较旧的 SiFive GCC-Metal 10.2.0 工具链，其链接器无法识别官方运行库目标文件中的 \texttt{zmmul} ISA 属性。检查 \texttt{type -a} 后发现课程工具链已经安装在 \texttt{/home/lenovo/riscv-elf-toolchains/bin}，将该目录置于 \texttt{PATH} 前端并重新生成 \texttt{main.o} 后，官方 \texttt{libsysy\_riscv.a} 可以正常链接。该现象说明预编译静态库中的目标架构属性需要与所使用的汇编器、链接器版本兼容。

\labfigure{figures/shared/q2_riscv/q2_1.png}{RV64 目标文件检查与 SysY 运行时库链接}{fig:q2-rv-link}

\subsubsection{运行结果验证}
链接成功后，\texttt{file} 显示 \texttt{main} 为 \texttt{ELF 64-bit LSB executable, UCB RISC-V}，并且采用静态链接。使用 \texttt{qemu-riscv64} 执行程序，对正常数据、包含负数的数据以及全奇数数据分别进行测试。运行时库额外打印的 \texttt{TOTAL} 计时行由 SysY 运行时库产生，不属于本程序的两个计算结果。

\begin{lstlisting}[language=bash]
qemu-riscv64 ./main
\end{lstlisting}
\labfigure{figures/shared/q2_riscv/q2_2.png}{可执行文件类型与第 1 组运行结果}{fig:q2-rv-test1}
\labfigure{figures/shared/q2_riscv/q2_3.png}{第 2 组运行结果（含负数）}{fig:q2-rv-test2}
\labfigure{figures/shared/q2_riscv/q2_4.png}{第 3 组运行结果（全奇数）}{fig:q2-rv-test3}

\begin{table}[H]
\centering
\caption{RISC-V 程序多组测试结果}
\begin{tabularx}{0.94\linewidth}{c X c c c}
\toprule
组别 & 输入 & 预期输出 & 实际输出 & 结果 \\
\midrule
1 & 5；1 2 -3 4 6 & 12；4 & 12；4 & 正确 \\
2 & 4；-2 3 8 -5 & 6；2 & 6；2 & 正确 \\
3 & 5；1 3 5 7 9 & 0；5 & 0；5 & 正确 \\
\bottomrule
\end{tabularx}
\end{table}

三组测试的实际输出均与对 SysY 源程序的直接计算结果一致。其中第 2 组覆盖了负数参与偶数求和的情况，第 3 组使偶数和为 0，可进一步验证条件分支和循环控制没有只针对单一输入工作。

\subsubsection{运行时库链接验证}
为了进一步确认最终可执行文件确实包含 SysY 运行时库，使用 \texttt{nm} 查看链接后的符号表：
\begin{lstlisting}[language=bash]
riscv64-unknown-elf-nm main \
  | grep -E "getint|putint|putch|sum_even|count_pos| main$"
\end{lstlisting}
\labfigure{figures/shared/q2_riscv/q2_5.png}{最终可执行文件中的自定义函数与 SysY 运行时库符号}{fig:q2-rv-nm}

符号表中同时存在手写汇编定义的 \texttt{sum\_even}、\texttt{count\_pos}、\texttt{main}，以及运行时库提供的 \texttt{getint}、\texttt{putint}、\texttt{putch}。由此可以确认最终程序不仅在运行结果上正确，而且确实完成了题目要求的“RISC-V 汇编程序链接 SysY 运行库”。

\subsubsection{RISC-V 部分小结}
本实验从一份包含数组、赋值、运算、条件分支、循环和函数等语言特性的 SysY 程序出发，手工完成了等价 RV64 汇编实现。数组元素访问被转换为地址计算和 \texttt{lw/sw}，\texttt{while} 与 \texttt{if} 被转换为条件分支和标签跳转，函数参数和返回值按照 RISC-V 调用约定通过 \texttt{a0}、\texttt{a1} 等寄存器传递。\texttt{main.s} 经汇编得到可重定位目标文件，再与官方 \texttt{libsysy\_riscv.a} 静态链接得到 RISC-V 可执行文件。三组输入在 QEMU 下均得到预期结果，\texttt{nm} 也验证了 SysY 运行时库函数已经进入最终可执行文件，说明本部分完成了从 SysY 语义到 RISC-V 目标程序的等价实现与运行验证。
